% !TEX TS-program = xelatex
% !TEX encoding = UTF-8 Unicode
% !Mode:: "TeX:UTF-8"

\documentclass{../../styles/resume}
\usepackage{../../styles/zh_CN-Adobefonts_external}
\usepackage{../../styles/linespacing_fix}
\usepackage{indentfirst}
\usepackage{cite}
\usepackage{setspace}
\onehalfspacing

\begin{document}

\pagenumbering{gobble}

\begin{center}
  \begin{minipage}{\textwidth}
    \centering
    \name{孟高翔}
    % \vspace{0.1em}
    \basicInfo{
      \email{1514190951@qq.com} \textperiodcentered\ 
      \phone{(+86) 136-7557-3771}
    }
    \basicInfo{
      {男 | 安徽泗县 | 中共党员}
    }
  \end{minipage}
\end{center}

\begin{tikzpicture}[remember picture, overlay] 
\node[anchor=north east, inner sep=0pt] at ([xshift=-2.5cm, yshift=-0.7cm]current page.north east) { 
\includegraphics[width=1in]{../../assets/images/photo.jpg} }; 
\end{tikzpicture}


%--------------------------------------
\section{\faGraduationCap\ 教育背景}
% \datedline{\textbf{武汉大学（双一流A类，985高校）|计算机学院|软件工程专业}}{\textit{硕士}\quad \textit{2026.09 -- 2028.06}}
\datedline{\textbf{武汉大学（双一流A类，985高校）|计算机学院|软件工程专业}}{\textit{本科}\quad \textit{2022.09 -- 2026.06}}
\begin{itemize}[leftmargin=*, labelsep=0.5em]
\item \textbf{本科绩点}：3.883\quad \textbf{专业排名}：16/218（7\%）
\item \textbf{荣誉奖项}：校级乙等奖学金（前30\%）、三好学生、优秀学生干部、甲等新生奖学金
\item \textbf{核心课程}：计算机网络（96），数据结构（92），操作系统（91），计算机组成与原理（90），软件需求与建模（90），云计算（92），线性代数（96）
\end{itemize}

%--------------------------------------
\section{\faUsers\ 实习经历}
\datedsubsection{\textbf{腾讯公司-} | 后台开发工程师}{2022.09 -- 2026.06}
\role{（腾讯控股有限公司）}{}
\begin{itemize}[leftmargin=*, labelsep=0.5em]
  \item \textbf{负责工作}：实习期间工作主要集中在xxx
  \item \textbf{工作收获}：……
\end{itemize}


% Section: Project Experience
\section{\faProjectDiagram\ 项目经历}

% Project 1: Star-SaaS Short Link
\datedsubsection{\textbf{星辰Star-SaaS短链接（个人项目）}}{2024.01 -- 2024.02}
\role{后端开发}{SpringBoot, SpringCloudAlibaba, RocketMQ, ShardingSphere, Redis, MySQL, Sentinel}
\begin{itemize}[leftmargin=*, labelsep=0.5em]
  \item \textbf{项目介绍}：SaaS短链接系统，提供高效、安全、可靠的短链接管理平台，简化长链接的管理与分享过程，并提供深入的分析与跟踪功能，助力用户灵活管理链接以提升营销效果。
  \item \textbf{线上地址}：\url{https://horiik} \quad \textbf{仓库地址}：\url{https://github.com}
  \item \textbf{功能描述}：
    \begin{itemize}
      \item 采用布隆过滤器快速判断短链接是否存在，性能远超分布式锁结合数据库查询的方案。
      \item 利用RocketMQ消息队列的削峰填谷特性，支持海量访问场景下的监控信息存储，实现存储与访问解耦。
      \item 通过更新数据库后删除缓存的策略，确保缓存与数据库之间的数据一致性。
      \item 基于Redis实现消息队列消费的幂等场景，保障消息在一定时间内仅被消费一次。
      \item 结合读写锁与RocketMQ延迟队列，支持短链接分组在高并发场景下的数据修改。
      \item 使用Sentinel实现接口QPS限流，触发限流规则后进行降级处理，确保系统稳定运行。
    \end{itemize}
  \item \textbf{项目收获}：通过开发本项目，深入理解并发场景下的开发设计，对SaaS架构、微服务、网关及流量控制有深刻实践经验，提升了分布式系统开发能力。
\end{itemize}
% Project 2: Star-RPC Framework
\datedsubsection{\textbf{星辰RPC框架（个人项目）}}{2024.03 -- 2024.04}
\role{系统开发}{Spring, Netty, Kryo, Zookeeper, Guava; C++11: ReactorTCP, Protobuf, RPC, TinyXML}
\begin{itemize}[leftmargin=*, labelsep=0.5em]
  \item \textbf{项目介绍}：星辰RPC框架，模拟实现了一个基础RPC框架，支持注册中心、网络传输、序列化、动态代理及负载均衡功能，同时提供传统BIO方式的Socket传输及JDK序列化作为对比方案。
  \item \textbf{仓库地址}：\url{https://github.com}
  \item \textbf{功能描述}：
    \begin{itemize}
      \item 使用Zookeeper作为注册中心，负责服务地址的注册与查找。
      \item 基于Netty实现RPC网络请求传输，包含调用的类名、方法名及参数。
      \item 采用Kryo序列化机制，优化性能并解决JDK序列化效率低及安全漏洞问题。
      \item 实现特定负载均衡策略，避免单服务器响应同一请求导致宕机或崩溃。
    \end{itemize}
  \item \textbf{项目收获}：通过开发本框架，掌握了RPC架构设计与实现，深化了对分布式系统、网络传输及服务治理的理解。
\end{itemize}
% Project 3: Axin-Star-Admin
\datedsubsection{\textbf{Axin-Star-Admin（个人项目）}}{2023.07 -- 2023.12}
\role{全栈开发}{SpringBoot, MyBatis, JWT, Redis, MySQL, Vue, TypeScript, Docker, Git}
\begin{itemize}[leftmargin=*, labelsep=0.5em]
  \item \textbf{项目介绍}：前后端分离的管理系统，基于RBAC模型实现细粒度权限控制，涵盖用户、角色、菜单管理，并扩展了Blog和Game管理模块，优化个人网页发布流程。
  \item \textbf{线上地址}：\url{https://admin} \quad \textbf{仓库地址}：\url{https://github.com}
  \item \textbf{功能描述}：
    \begin{itemize}
      \item 基于RBAC模型实现接口及按钮级别的细粒度权限控制，支持用户、角色、菜单管理。
      \item Blog模块支持文章、分类、标签管理，简化网页发布流程，定时任务存储草稿避免丢失，图片上传至阿里云OSS图床。
      \item 支持一键发布文章至Hexo前台目录，调用shell脚本重新部署并刷新CDN。
      \item 采用阿里云STS临时凭证访问OSS及CDN，凭证缓存至内存，使用线程池定时清理过期缓存。
      \item 编写Dockerfile制作Java后端镜像，结合docker-compose集成MySQL、Redis、Nginx及OpenJDK镜像，实现一键部署所有容器。
    \end{itemize}
  \item \textbf{项目收获}：通过开发本系统，优化了项目统一管理能力，深入实践后端主流技术栈及前后端分离开发流程，提升了全栈开发与容器化部署经验。
\end{itemize}
%--------------------------------------
% 定义科研与竞赛经历部分
\section{\faTrophy\ 科研与竞赛经历}
\datedline{\textit{省级大创—智能招聘信息匹配平台 | 核心成员}}{2025.03 -- 2025.06}
\datedline{\textit{2024年全国大学生数学建模竞赛湖北赛区三等奖 | 核心成员}}{2024.09}
\datedline{\textit{第九届华为ICT大赛编程赛省赛一等奖 | 核心成员}}{2025.04}
\datedline{\textit{第十五届蓝桥杯湖北赛区二等奖 | 个人比赛}}{2024.04}

%--------------------------------------
\section{\faInfo\ 其他信息}
\begin{itemize}
  \item 英语水平：雅思6.5、CET-4、CET-6
  \item 专业技能：熟悉 Java、C++、Go、TypeScript、Python；熟悉常见前后端框架；熟悉 MySQL、Linux、Git、Maven、Docker、K8s 等
  \item 兴趣爱好：体育竞赛（院运会篮球赛铜牌、排球赛第五名）、篮球、游泳
\end{itemize}

\end{document}
